\section{}
\setcounter{subsection}{2}
\subsection{Desarrollo}
A continuación se detallan las actividades ejecutadas durante el período de estadías, organizadas por semanas. Estas tareas se enfocaron en la \textbf{optimización de infraestructura tecnológica} de la institución, con el objetivo de mejorar la eficiencia operativa, garantizar la seguridad de los sistemas y cumplir con los estándares de calidad bajo la norma \textbf{ISO 9001}.  

Cada semana se abordaron proyectos críticos, como actualizaciones de sistemas, mantenimiento preventivo y análisis de red, que contribuyeron directamente a:  
\begin{itemize}
	\item \textbf{Reducir vulnerabilidades técnicas} (ej. migración a Windows 11, actualizaciones de seguridad).
	\item \textbf{Optimizar recursos} (ej. reorganización física de laboratorios para cursos especializados).
	\item \textbf{Fortalecer procesos institucionales}.
\end{itemize}

\subsubsection{Tareas por semana:}
\subsubsubsection{Tercera Semana de Febrero 2025}
\begin{itemize}
	\item \textbf{Actualización de infraestructura}:
	\begin{itemize}
		\item Migración de 19/24 equipos del laboratorio a Windows 11
		\item Identificación y diagnóstico de 5 equipos con problemas de actualización (falta de periféricos, credenciales administrativas y fallas de BIOS)
		\item Reparación de cable Ethernet en sala de conferencias con recableado según estándar T568A
	\end{itemize}
	
	\item \textbf{Gestión de dispositivos móviles}:
	\begin{itemize}
		\item Evaluación de 15 laptops UMA: 13 actualizadas, 2 marcadas para retiro (UMA003 y UMA006 por obsolescencia y falta de licencias)
		\item Implementación de solución temporal para etiquetado duplicado (UMA015-b)
	\end{itemize}
	
	\item \textbf{Soporte técnico especializado}:
	\begin{itemize}
		\item Configuración de impresora Lanier para formato oficio
		\item Diagnóstico de teléfono IP con falla en puerto PoE (oficina DAP)
		\item Reparación de altavoz Bluetooth en sala de videoconferencias
	\end{itemize}
\end{itemize}

\subsubsubsection{Cuarta Semana de Febrero 2025}
\begin{itemize}
	\item \textbf{Seguimiento de pendientes}:
	\begin{itemize}
		\item Actualización exitosa de 3 equipos pendientes mediante coordinación con sysadmin
		\item Retiro definitivo de laptops UMA por decisión institucional
	\end{itemize}
	
	\item \textbf{Optimización de red}:
	\begin{itemize}
		\item Recableado y recrimpado de 12 patch cords defectuosos
		\item Actualización masiva de firmware y controladores en 7 áreas clave usando Dell Command Update
	\end{itemize}
	
	\item \textbf{Incidentes críticos}:
	\begin{itemize}
		\item Resolución de problema DHCP en teléfono IP mediante reubicación de puerto en patch panel
		\item Diagnóstico de error de actualización en equipo de Coordinación (Elizabeth Persona)
	\end{itemize}
\end{itemize}

\subsubsubsection{Primera Semana de Marzo 2025}
\begin{itemize}
	\item \textbf{Mantenimiento preventivo}:
	\begin{itemize}
		\item Reemplazo de batería CMOS en equipos con código de error 3-1
		\item Limpieza interna de equipos afectados por tormentas de arena
	\end{itemize}
	
	\item \textbf{Gestión de incidencias}:
	\begin{itemize}
		\item Reparación emergente de puerto RJ45 en oficina de Vinculación con herramientas improvisadas
		\item Diagnóstico de falla en lector de huellas del Aula 1 mediante pruebas de control de acceso
	\end{itemize}
	
	\item \textbf{Proyectos especiales}:
	\begin{itemize}
		\item Análisis inicial de infraestructura de red para certificación ISO 9001
		\item Pruebas de conectividad para videollamadas con equipo de TI en Ciudad de México
	\end{itemize}
\end{itemize}

\textbf{Metodologías y Herramientas Utilizadas}
\begin{itemize}
	\item \textbf{Herramientas técnicas}:
	\begin{itemize}
		\item Dell Command Update para actualizaciones de firmware y BIOS
		\item Wireshark para análisis básico de tráfico de red
		\item PowerShell para automatización de despliegues de SO
	\end{itemize}
	
	\item \textbf{Protocolos de seguridad}:
	\begin{itemize}
		\item Estándar T568A para crimpado de cables
		\item NIST SP 800-88 para eliminación segura de datos
	\end{itemize}
\end{itemize}

\subsubsubsection{Cuarta Semana de Marzo 2025}
\begin{itemize}
	\item \textbf{Seguimiento de pendientes}:
	\begin{itemize}
		\item Punto de acceso del segundo piso sigue pendiente por falta de respuesta del proveedor en Chihuahua.
	\end{itemize}
	
	\item \textbf{Actualización de equipos}:
	\begin{itemize}
		\item Migración de 24 equipos del Laboratorio de Planta Baja de Windows 10 a Windows 11.
		\item Instalación de actualizaciones de seguridad y controladores.
		\item Reconfiguración física del espacio para optimizar el curso de LinkedIn.
		\item Limpieza de historiales de navegación post-curso.
	\end{itemize}
	
	\item \textbf{Equipo Portátil de Vinculación}:
	\begin{itemize}
		\item Migración exitosa a Windows 11 e instalación de Microsoft Office Suite.
		\item Diagnóstico de disco duro con alerta de falla (Dell/SMART/CrystalDisk), monitoreo continuo recomendado.
	\end{itemize}
	
	\item \textbf{Etiquetado de equipos}:
	\begin{itemize}
		\item Impresión y corte de 32 páginas de etiquetas.
		\item Entrega a Contador y Subdirección Administrativa para colocación física.
		\item Implementación de sistema de numeración para inventario.
	\end{itemize}
	
	\item \textbf{Mantenimiento de red}:
	\begin{itemize}
		\item Reconección física completa del punto de acceso del segundo piso (puertos y cableado).
		\item Pruebas de conectividad fallidas; incidencia reportado al ingeniero Benjamín en la Ciudad de Chihuahua
	\end{itemize}
	
	\item \textbf{Pendientes}:
	\begin{itemize}
		\item Intervención del equipo de red en Chihuahua para el punto de acceso del segundo piso.
		\item Verificación de colocación de etiquetas por parte de Contador.
		\item Entrega del Análisis de Red a más tardar el miércoles.
	\end{itemize}
\end{itemize}

\textbf{Metodologías y Herramientas Utilizadas}
\begin{itemize}
	\item \textbf{Herramientas técnicas}:
	\begin{itemize}
		\item Herramientas Dell para diagnóstico de hardware.
		\item SMART y CrystalDisk para análisis de disco duro.
		\item Microsoft Windows 11 para migración de sistemas.
	\end{itemize}
	
	\item \textbf{Protocolos de gestión}:
	\begin{itemize}
		\item Sistema de tickets para registro de incidencias.
		\item Estándares de inventario para etiquetado de equipos.
	\end{itemize}
\end{itemize}

\subsubsubsection{Primera Semana de Abril 2025}
\begin{itemize}
	\item \textbf{Seguimiento de pendientes}:
	\begin{itemize}
		\item No hay pendientes reportados para esta semana.
	\end{itemize}
	
	\item \textbf{Actualización de equipos}:
	\begin{itemize}
		\item Actualización de todos los equipos UMA.
		\item Verificación de cables y conexiones para garantizar su funcionamiento.
	\end{itemize}
	
	\item \textbf{Configuración de red y conectividad}:
	\begin{itemize}
		\item Configuración de un módem doméstico para permitir la conexión directa de los estudiantes al módem, evitando la conexión a través del teléfono del profesor en la Unidad Móvil de Aprendizaje (UMA).
		\item Finalización del diagnóstico de red que se llevó a cabo durante las últimas semanas en CIITA.
	\end{itemize}
	
	\item \textbf{Mantenimiento de hardware}:
	\begin{itemize}
		\item Diagnóstico y solución del problema con el teclado y mouse Logitech en la sala de conferencias. Se identificó que el adaptador USB estaba defectuoso y se reemplazó con otro modelo funcional de la misma marca.
	\end{itemize}
	
	\item \textbf{Mantenimiento de equipos periféricos}:
	\begin{itemize}
		\item Reemplazo del tóner en la copiadora Lanier.
	\end{itemize}
\end{itemize}

\textbf{Metodologías y Herramientas Utilizadas}
\begin{itemize}
	\item \textbf{Herramientas técnicas}:
	\begin{itemize}
		\item Diagnóstico de hardware para identificar fallas en dispositivos USB.
		\item Configuración de módems y redes domésticas.
	\end{itemize}
	

\subsubsection{Análisis de Red}
El análisis de red realizado en CIITA se centró en evaluar la infraestructura tecnológica existente, identificando vulnerabilidades críticas y proponiendo soluciones alineadas con estándares de seguridad y normativas institucionales (ISO 9001, Ley Orgánica del IPN). A continuación, se detallan los aspectos clave:

\paragraph{Metodologías y Herramientas}
\begin{itemize}
	\item \textbf{Diagnóstico de Segmentación}: 
	\begin{itemize}
		\item Uso de herramientas como \textit{Nmap} para analizar tráfico de red y topología.
		\item Identificación de problemas en la división con VLSM (Sección 3.1), incluyendo agotamiento de direcciones IP y falta de aislamiento entre subredes.
	\end{itemize}
	
	\item \textbf{Evaluación de Seguridad}:
	\begin{itemize}
		\item Auditoría del firewall \textit{FortiGate} (Sección 6.5.1) para detectar reglas obsoletas o permisivas.
		\item Análisis de riesgos en SSID (Sección 6.5.2–6.5.3), destacando la propagación lateral de amenazas en la red de invitados.
	\end{itemize}
	
	\item \textbf{Pruebas de Vulnerabilidad}:
	\begin{itemize}
		\item Escaneo con \textit{Fing} para detectar equipos terminales.
		\item Verificación de cumplimiento de políticas técnicas (Sección 6.7), como cifrado de dispositivos móviles y uso de Active Directory.
	\end{itemize}
\end{itemize}

\paragraph{Hallazgos Principales}
\begin{itemize}
	\item \textbf{Conectividad y Segmentación} (Sección 6.1–6.2):
	\begin{itemize}
		\item Diseño ineficiente con switches no configurados (ej. inundación CAM, Sección 6.4.1).
		\item Escasez de direcciones IP debido a asignación estática no optimizada.
	\end{itemize}
	
	\item \textbf{Seguridad} (Sección 6.5–6.9):
	\begin{itemize}
		\item Exposición del sistema HVAC a la red principal (Sección 6.6.6), violando normas ISO 9001.
		\item Falta de autenticación multifactor (2FA) y políticas de actualización (Sección 6.8).
	\end{itemize}
\end{itemize}

\paragraph{Conclusiones del Análisis}
Los resultados evidenciaron la necesidad de un rediseño de red (Sección 7.1) y la implementación de estrategias de defensa en profundidad (Sección 6.9). El reporte técnico completo, incluyendo datos detallados y evidencia, se adjunta en el \textbf{Anexo A} de este documento.